\documentclass[12pt]{article}
\usepackage[utf8]{inputenc}
\usepackage[T1]{fontenc}
\usepackage{lmodern}
\usepackage{microtype}
\usepackage[hidelinks]{hyperref}
\usepackage[a4paper, margin=3cm]{geometry}
\usepackage{parskip}

\title{How to get into Meditation}
\date{28/06/26}

\begin{document}
\maketitle

% [ADD: Write one arguable thesis sentence before drafting further.]

% ──────────────────────────────────────────────────────────────
% SECTION I: OPENING
% ──────────────────────────────────────────────────────────────

One of the most beautiful things in Sahaja Yoga is the bliss of the collective.
That is also what motivates us to travel to Cabella and other places of gathering.
Soak in all the good vibes, cleanse ourselves, get into a clarity, and then take it all home.

Once at home, everything feels different.
We're inspired and confident, and feel like we can take on the world and change it for the better.
That usually lasts for a week or two until daily life takes hold of us and we're back to the struggle of combining the idealism of Sahaj life with the daily reality of living in the normal world.
The main reason for this is that we are not able to maintain the depth in our meditation once the ease of Cabella-life is gone, where we're not confronted with our usual problems.
And we can change this.

% ──────────────────────────────────────────────────────────────
% SECTION II: THE DISCOVERY
% ──────────────────────────────────────────────────────────────

I know that situation very well.
The mystery of how to maintain a steady vibrational state --- just like in Cabella, but at home.
True meditation seemed like a matter of chance.
No matter how often I would balance my channels, how many Mantras I said, or even how many speeches I would listen to.
Sometimes I got thoughtlessness, and sometimes not.

Until I came across a guided meditation by Shri Mataji from 1979, titled \textit{How to get into meditation}.
Upon following her instructions I was completely struck and shook --- I have just learned how to actually meditate.
After 5 years of seriously practicing SY and a whole lifetime growing up in the SY community, it was a random and not particularly popular video of around 15 minutes that would be the key for my spiritual journey.

In the guided meditation the sequence is as follows.
First, the attention goes to the heart, where the lotus feet of the Guru are placed in the heart.
Then, there is space to balance the channels or invoke Shri Ganesha to support the rising of the Kundalini.
After that the attention goes towards the mind with the goal to steady the attention.
For that the question: ``Where is my attention?'' gives guidance.
Once the attention is stabilized and held in control, the \textit{Nirvichara} Mantra is chanted to create the complete cessation of thought, i.e.\ thoughtlessness.
The opening of Agya follows only after that, namely by bringing down the Ego.
Either with the \textit{Mahat Ahankara} Mantra, or literally by bringing the Ego down using the right hand that lifts the left side and lowers the right.
Now, the real meditation starts as the Kundalini passes through the Agya and touches the Sahastrara, enlightening the being with the light of the spirit and establishing a connection, or union, with the collective unconscious.

For years until now this had become the method for me to get into meditation on a daily basis.
Always the same.
Morning and evening.
And I must say, I don't understand why almost no guided meditations in Sahaja Yoga follow this procedure, and neither did I understand the theory behind the method.

% ──────────────────────────────────────────────────────────────
% SECTION III: THE FRAMEWORK
% ──────────────────────────────────────────────────────────────

That changed this year in January, when we attended the Brazil Tour and the subsequent Shivaratri seminar in Brasilia.
In the inauguration, or the welcome ceremony, of the seminar, Uncle Eduardo Marino gave a presentation on attention and how to meditate.
He would present his understanding of attention and its importance for the meditation compiled from different speeches of SM.
Additionally, he would take Patanjali's \textit{Yoga Sutras} and the 8 limbs of Yoga to support the message, and lastly he guided a meditation that would put in practice the presented theory.
For me it was an absolute excitement, as I had finally attended a Sahaj event, where the guided meditation resembled very well the procedure from that 1979 guided meditation by Shri Mataji.
Later, I also started understanding the reason why this method worked so well.

% [OPTIONAL: Consider including Uncle Eduardo's definition of attention before presenting Patanjali:
%  "the attention is the divine within us --- the spectator of a screen that our mind uses to project its
%  images on. Depending on the state of our awareness, our attention gets dragged towards these projections,
%  which creates a thought process and takes us into the realm of the mind, i.e. the realm of illusion."
%  Decision: include or skip. Adds depth but may slow the pace here.]

As Uncle Eduardo would point out in his presentation there are 8 limbs in the \textit{Yoga Sutras}, where the first 4 represent the outward practice, like \textit{Pranayama}, or \textit{Asana}, while the final 4 represent the inward practice.

% [NOTE: The four stages appear below as a numbered list AND as prose sentences — remove one.
%  List: cleaner visual reference. Prose: Samadhi definition has more room to breathe.]

1.\ \textit{Pratyahara} (Moving the Attention Inwards) \\
2.\ \textit{Dharana} (Concentrating the Attention) \\
3.\ \textit{Dhyana} (Ceasing Thoughts) \\
4.\ \textit{Samadhi} (Thoughtless Awareness)

\textit{Pratyahara} means the drawing of the attention towards the inside.
\textit{Dharana} describes the focusing of the attention, and \textit{Dhyana} the complete cessation of thought.
And finally, \textit{Samadhi} means the state of awareness, where Yoga happens, or when we are in Yoga, i.e.\ union.

% ──────────────────────────────────────────────────────────────
% SECTION IV: THE MAPPING
% ──────────────────────────────────────────────────────────────

Shri Mataji's guided meditation connects to those steps in the following way.

Firstly, the attention on the heart, with the purpose of inviting the Guru into our heart --- that is the \textit{Pratyahara}.
We pull our attention inward completely and look for the divinity within.
After that, the question: ``Where is my attention?'' represents the \textit{Dharana}.
Here we collect and focus our attention that is in movement due to mental projections from the mind that drag our attention towards it.
Once the attention is steady the \textit{Dhyana} begins that consists of increasing the gap between the thoughts until a complete stillness is achieved.
For this Shri Mataji recommends the \textit{Nirvichara} Mantra.
And finally, the \textit{Mahat Ahankara} Mantra represents the beginning of the final phase --- the \textit{Samadhi} --- where a complete oneness and connection with the divine is achieved by opening the Sahastrara.

\begin{tabular}{lll}
\textit{Pratyahara} & Moving the attention inward & Taking the attention to the heart            \\
\textit{Dharana}    & Concentrating the attention & Watching the thoughts                        \\
\textit{Dhyana}     & Ceasing thoughts            & Increasing the gap between thoughts          \\
\textit{Samadhi}    & Thoughtless awareness       & Opening Agya and entering the kingdom of God \\
\end{tabular}

% ──────────────────────────────────────────────────────────────
% SECTION V: THE KEY INSIGHT
% ──────────────────────────────────────────────────────────────

One question I often had was what role the Kundalini played in this process.
Do I need an awakened Kundalini to be thoughtless?
This is at least something that Sahaja Yogis often answer with: `yes.'
If that were true we would find that only Kundalini awakening leads to measurable thoughtlessness via EEGs, but that is not the case.
Buddhists and practitioners of other types of meditation achieve measurable stillness in their brainwaves just as well, and we know that the vast majority didn't receive their Self-Realization, or Kundalini awakening.

And therefore, when we talk of thoughtless awareness in Sahaja Yoga we actually mean \textit{Samadhi}.
That is a very important distinction because only for this final stage of Yoga the awakening of the Kundalini is mandatory.
All the previous steps, from \textit{Pratyahara} to \textit{Dhyana}, don't actually require Kundalini awakening.

% [NOTE: Consider adding this sentence to close the insight:
%  "Where the Kundalini becomes necessary is for passing through the Agya ---
%  or for entering Samadhi in Patanjali's words."]

% ──────────────────────────────────────────────────────────────
% SECTION VI: THE CHALLENGE
% ──────────────────────────────────────────────────────────────

Until \textit{Dhyana}, it's mostly in our hands.
To turn the attention inward, and focus it, is absolutely not Kundalini's job.
It is the self-respect we have to show to make ourselves worthy for our Kundalini to rise in the first place.
From \textit{Dhyana} to \textit{Samadhi}, on the other hand, we are completely left to her will to bring us anywhere.
Opening the Agya cannot be achieved by our individual effort nor by our spiritual prowess.
Only the grace of the Kundalini allows the passage through the narrow gate and towards Yoga.

Therefore, getting into meditation should not be a stochastic event, that depending on the weather, or the alignment of the stars, it may happen to us or not.
On the contrary it should be the strongest tool of a Sahaja Yogi, a reliable process to align ourselves with the Paramchaitanya --- whenever we desire to.
For that we just need more method.
We need to find out what disturbs our attention, where its movement goes and how to control it.
We also need to practice focus and honesty with ourselves.
It should be transparent to us what desires lurk inside and how our Ego tries to escape its exposure.

% [ADD: A sentence that directly names the accountability — that this is not someone else's job to teach you.
%  Consider: "In other words, it's our responsibility to find out how we can keep up the strong vibrational
%  state after our Cabella trip." — this circles back to the opening and creates a satisfying structure.]

% ──────────────────────────────────────────────────────────────
% SECTION VII: THE ANALOGY
% ──────────────────────────────────────────────────────────────

% [CONDENSE: This paragraph runs too long. Target: the fall and the mother beginning treatment
%  in no more than 3 sentences total before the counter-scenario.]
Say, we are a small child and there is our mother that is taking care of us.
She represents our Kundalini, who protects us, nurtures us, and looks after us.
Imagine the child playing outside, running around, jumping over obstacles, and being fully immersed in its imagination --- just like us how we manoeuvre through life.
Eventually, something goes wrong and the child falls badly and hurts itself.
The child runs to its mother, crying loudly, and with great concern the mother starts tending to the child.
First, calming it down, then sitting it down to look at the wound.
After a quick diagnosis she will then prepare her treatment that consists of disinfecting the wound and applying a bandage to stop the bleeding.

Ideally, the child would stay calm and the mother slowly and carefully executes all parts of the treatment with great care for the wounded knee.
However, now imagine the child would continue running around in the house.
How would the mother be able to help the child?
She tries to run behind the child, but doesn't manage to apply the disinfectant properly.
Further, the bandage she could manage to put on wouldn't hold properly because of the constant movement of the child.

Now, think of the child as our mind.
How can the Kundalini tend to our problems if we cannot steady our attention?
She cannot.
Not because of her inability, but because of her respect of our free will.
We have to be mature enough to prepare ourselves accordingly so that she can treat us with all the love, affection and care that she wishes to.
That is the reason why \textit{Pratyahara}, \textit{Dharana} and \textit{Dhyana} are for us to provide to our Kundalini, so that she can take us to \textit{Samadhi}, to the Sahastrara.

% [NOTE: Consider adding this as a closing sentence of the analogy:
%  "Just like the boy better just watch the process instead of asking too many questions about the plaster,
%  or giving advice on how to apply the disinfectant properly."]

% ──────────────────────────────────────────────────────────────
% SECTION VIII: CLOSE
% ──────────────────────────────────────────────────────────────

There is a lot we can do to improve our meditation, especially when it comes to calming our mind.
We don't have to leave it to chance, the weather, or the current planetary constellation whether we go thoughtless or not.
We can simply practice the control over our mind, so that our Kundalini can perform her miracle.
The method described here, inspired by Shri Mataji's guided meditation, may serve as inspiration, but ultimately it's up to every individual to find out their most reliable method.
Let's find out how to keep the magic alive after Cabella.

\section*{Resources}

\begin{itemize}
  \item \href{https://drive.proton.me/urls/CB641G50BR#RgvXq2i1MiLc}{Uncle Eduardo's presentation on attention}
  \item \href{https://www.youtube.com/watch?v=RBUCsIEYK_Y}{Shri Mataji's guided meditation (1979)}
\end{itemize}

\end{document}
